% GENERATED FILE. Edit canonical lessons, then run npm run generate:books.
% canonical-source-hash: fnv1a64:1ad09ae76b357198
% canonical-lessons: TA-C92-slowly, TA-C92-again, TA-C92-loudly, TA-C92-oneminute, TA-C92-wait

\chapter{Slowly, Again, Louder}
\label{ch:ta-slowly-again-louder}
\hlchaptermodality{92}

\begin{chapteropening}
\textbf{By the end of this chapter:} \emph{I can ask someone to speak slowly, say it again, speak louder, and wait a minute.}

Complete the last lesson of chapter 92: i can ask someone to speak slowly, say it again, speak louder, and wait a minute.
\end{chapteropening}

\section[\texorpdfstring{\ta{மெதுவாக} (metuvāka)}{மெதுவாக (metuvāka)}]{\ta{மெதுவாக} (metuvāka) — slowly}

\label{lesson:TA-C92-slowly}



\begin{quote}
Before the new one: say the Tamil for I am sleepy, then the Tamil for well.
\end{quote}

\subsection*{You'll want to know: \ta{மெதுவாக}}
\textbf{\ta{மெதுவாக}} (\emph{metuvāka}) — \textquotedblleft{}slowly\textquotedblright{}.

\textbf{\ta{மெதுவாகப்} \ta{பேசுங்கள்}} — \textquotedblleft{}please speak slowly\textquotedblright{}, with the verb \textbf{\ta{பேசு}} you read long ago.

\textbf{In the family, and next door}

\noindent
\begin{tabularx}{\linewidth}{@{}>{\raggedright\arraybackslash}X>{\raggedright\arraybackslash}X@{}}
\toprule
\textbf{} & \textbf{word} \\
\midrule
Kannada & \ka{ನಿಧಾನವಾಗಿ} (\emph{nidhānavāgi}) \\
Telugu & \te{నెమ్మదిగా} (\emph{nemmadigā}) \\
Malayalam & \ml{പതുക്കെ} (\emph{patukke}) \\
Hindi (neighbour) & \dv{धीरे} (\emph{dhīre}) \\
English & slowly \\
\bottomrule
\end{tabularx}

\begin{grammarlens}[title={What you've built}]
The repair a beginner needs most.
\end{grammarlens}

\subsection*{Your turn}
\begin{itemize}
\raggedright
  \item \emph{Say it:} \emph{metuvāka}
  \item \emph{Say it:} \emph{metuvāka}, once more
  \item \emph{Say it:} say \emph{metuvākap pēcuṅkaḷ}
  \item \emph{Recall:} say the Tamil for I am sleepy, then the Tamil for well, then say \emph{metuvāka} again
\end{itemize}

\begin{tcolorbox}[breakable,colback=teal!4,colframe=teal!35!black,title={Before you move on}]
What does \ta{மெதுவாக} mean? (\textquotedblleft{}Slowly\textquotedblright{}.) Say it once more.
\end{tcolorbox}
\section[\texorpdfstring{\ta{மறுபடியும்} (maṟupaṭiyum)}{மறுபடியும் (maṟupaṭiyum)}]{\ta{மறுபடியும்} (maṟupaṭiyum) — again}

\label{lesson:TA-C92-again}



\begin{quote}
Before the new one: say the Tamil for well, then the Tamil for slowly.
\end{quote}

\subsection*{You'll want to know: \ta{மறுபடியும்}}
\textbf{\ta{மறுபடியும்}} (\emph{maṟupaṭiyum}) — \textquotedblleft{}again\textquotedblright{}.

\textbf{\ta{மறுபடியும்} \ta{சொல்லுங்கள்}} — \textquotedblleft{}please say it again\textquotedblright{}.

\textbf{In the family, and next door}

\noindent
\begin{tabularx}{\linewidth}{@{}>{\raggedright\arraybackslash}X>{\raggedright\arraybackslash}X@{}}
\toprule
\textbf{} & \textbf{word} \\
\midrule
Hindi (neighbour) & \dv{दोबारा} (\emph{dobārā}) \\
English & again \\
\bottomrule
\end{tabularx}

\begin{grammarlens}[title={What you've built}]
Slowly, and again.
\end{grammarlens}

\subsection*{Your turn}
\begin{itemize}
\raggedright
  \item \emph{Say it:} \emph{maṟupaṭiyum}
  \item \emph{Say it:} \emph{maṟupaṭiyum}, once more
  \item \emph{Say it:} say \emph{maṟupaṭiyum colluṅkaḷ}
  \item \emph{Recall:} say the Tamil for well, then the Tamil for slowly, then say \emph{maṟupaṭiyum} again
\end{itemize}

\begin{tcolorbox}[breakable,colback=teal!4,colframe=teal!35!black,title={Before you move on}]
What does \ta{மறுபடியும்} mean? (\textquotedblleft{}Again\textquotedblright{}.) Say it once more.
\end{tcolorbox}
\section[\texorpdfstring{\ta{சத்தமாக} (cattamāka)}{சத்தமாக (cattamāka)}]{\ta{சத்தமாக} (cattamāka) — loudly}

\label{lesson:TA-C92-loudly}



\begin{quote}
Before the new one: say the Tamil for slowly, then the Tamil for again.
\end{quote}

\subsection*{You'll want to know: \ta{சத்தமாக}}
\textbf{\ta{சத்தமாக}} (\emph{cattamāka}) — \textquotedblleft{}loudly\textquotedblright{}.

\textbf{\ta{சத்தம்}} is \textquotedblleft{}a noise, a sound\textquotedblright{}. \textbf{\ta{கொஞ்சம்} \ta{சத்தமாகப்} \ta{பேசுங்கள்}} — \textquotedblleft{}please speak a little louder\textquotedblright{}.

\textbf{In the family, and next door}

\noindent
\begin{tabularx}{\linewidth}{@{}>{\raggedright\arraybackslash}X>{\raggedright\arraybackslash}X@{}}
\toprule
\textbf{} & \textbf{word} \\
\midrule
Kannada & \ka{ಜೋರಾಗಿ} (\emph{jōrāgi}) \\
Telugu & \te{గట్టిగా} (\emph{gaṭṭigā}) \\
Malayalam & \ml{ഉറക്കെ} (\emph{uṟakke}) \\
Hindi (neighbour) & \dv{ज़ोर} \dv{से} (\emph{zor se}) \\
English & loudly \\
\bottomrule
\end{tabularx}

\begin{grammarlens}[title={What you've built}]
Slowly, again, louder.
\end{grammarlens}

\subsection*{Your turn}
\begin{itemize}
\raggedright
  \item \emph{Say it:} \emph{cattamāka}
  \item \emph{Say it:} \emph{cattamāka}, once more
  \item \emph{Say it:} say \emph{koñcam cattamākap pēcuṅkaḷ}
  \item \emph{Recall:} say the Tamil for slowly, then the Tamil for again, then say \emph{cattamāka} again
\end{itemize}

\begin{tcolorbox}[breakable,colback=teal!4,colframe=teal!35!black,title={Before you move on}]
What does \ta{சத்தமாக} mean? (\textquotedblleft{}Loudly\textquotedblright{}.) Say it once more.
\end{tcolorbox}
\section[\texorpdfstring{\ta{ஒரு} \ta{நிமிடம்} (oru nimiṭam)}{ஒரு நிமிடம் (oru nimiṭam)}]{\ta{ஒரு} \ta{நிமிடம்} (oru nimiṭam) — one minute, a moment}

\label{lesson:TA-C92-oneminute}



\begin{quote}
Before the new one: say the Tamil for again, then the Tamil for loudly.
\end{quote}

\subsection*{You'll want to know: \ta{ஒரு} \ta{நிமிடம்}}
\textbf{\ta{ஒரு} \ta{நிமிடம்}} (\emph{oru nimiṭam}) — \textquotedblleft{}one minute\textquotedblright{}, said the way English says \textquotedblleft{}a moment\textquotedblright{}.

\textbf{\ta{ஒரு}} is \textquotedblleft{}one\textquotedblright{}; \textbf{\ta{நிமிடம்}} is \textquotedblleft{}a minute\textquotedblright{}.

\textbf{In the family, and next door}

\noindent
\begin{tabularx}{\linewidth}{@{}>{\raggedright\arraybackslash}X>{\raggedright\arraybackslash}X>{\raggedright\arraybackslash}X@{}}
\toprule
\textbf{} & \textbf{word} & \textbf{same root} \\
\midrule
Kannada & \ka{ಒಂದು} \ka{ನಿಮಿಷ} (\emph{ondu nimiṣa}) & yes \\
Telugu & \te{ఒక్క} \te{నిమిషం} (\emph{okka nimiṣaṁ}) & yes \\
Hindi (neighbour) & \dv{एक} \dv{मिनट} (\emph{ek minaṭ}) &  \\
English & one minute, a moment &  \\
\bottomrule
\end{tabularx}

\begin{grammarlens}[title={What you've built}]
A way to buy time.
\end{grammarlens}

\subsection*{Your turn}
\begin{itemize}
\raggedright
  \item \emph{Say it:} \emph{oru nimiṭam}
  \item \emph{Say it:} \emph{oru nimiṭam}, once more
  \item \emph{Say it:} say \emph{oru nimiṭam}
  \item \emph{Recall:} say the Tamil for again, then the Tamil for loudly, then say \emph{oru nimiṭam} again
\end{itemize}

\begin{tcolorbox}[breakable,colback=teal!4,colframe=teal!35!black,title={Before you move on}]
What does \ta{ஒரு} \ta{நிமிடம்} mean? (\textquotedblleft{}One minute, a moment\textquotedblright{}.) Say it once more.
\end{tcolorbox}
\section[\texorpdfstring{\ta{காத்திருங்கள்} (kāttiruṅkaḷ)}{காத்திருங்கள் (kāttiruṅkaḷ)}]{\ta{காத்திருங்கள்} (kāttiruṅkaḷ) — please wait}

\label{lesson:TA-C92-wait}



\begin{quote}
Before the new one: say the Tamil for loudly, then the Tamil for one minute.
\end{quote}

\subsection*{You'll want to know: \ta{காத்திருங்கள்}}
\textbf{\ta{காத்திருங்கள்}} (\emph{kāttiruṅkaḷ}) — \textquotedblleft{}please wait\textquotedblright{}.

\textbf{\ta{காத்திருங்கள்}, \ta{ஒரு} \ta{நிமிடம்}} — \textquotedblleft{}please wait, one minute\textquotedblright{}.

\textbf{In the family, and next door}

\noindent
\begin{tabularx}{\linewidth}{@{}>{\raggedright\arraybackslash}X>{\raggedright\arraybackslash}X>{\raggedright\arraybackslash}X@{}}
\toprule
\textbf{} & \textbf{word} & \textbf{same root} \\
\midrule
Kannada & \ka{ಕಾಯಿರಿ} (\emph{kāyiri}) & yes \\
Telugu & \te{ఆగండి} (\emph{āgaṇḍi}) &  \\
Malayalam & \ml{കാത്തിരിക്കൂ} (\emph{kāttirikkū}) & yes \\
Hindi (neighbour) & \dv{रुकिए} (\emph{rukie}) &  \\
English & please wait &  \\
\bottomrule
\end{tabularx}

\begin{grammarlens}[title={What you've built}]
Slowly, again, louder, one minute, wait.
\end{grammarlens}

\subsection*{Your turn}
\begin{itemize}
\raggedright
  \item \emph{Say it:} \emph{kāttiruṅkaḷ}
  \item \emph{Say it:} \emph{kāttiruṅkaḷ}, once more
  \item \emph{Say it:} say \emph{kāttiruṅkaḷ, oru nimiṭam}
  \item \emph{Recall:} say the Tamil for loudly, then the Tamil for one minute, then say \emph{kāttiruṅkaḷ} again
\end{itemize}

\begin{tcolorbox}[breakable,colback=teal!4,colframe=teal!35!black,title={Before you move on}]
What does \ta{காத்திருங்கள்} mean? (\textquotedblleft{}Please wait\textquotedblright{}.) Say it once more.
\end{tcolorbox}
